\subsection{Statistical analysis}
\label{sec:stats}

We estimate training variation separately from evaluation variation.
\emph{Training variation} is estimated by repeating every colour-augmentation
configuration at seeds 0, 1, and 2 under an identical protocol. We also repeat
the \texttt{stock} configuration so that its variation across runs provides a
noise floor for differences between recipes. Results are reported as the mean
over seeds with the observed range. We do not assume that three runs establish
a normal distribution.

\emph{Evaluation variation} is estimated at image level, where the sample is
larger. For two models at one perturbation level, we calculate their difference
in counting error for each of the 161 validation images. We resample these
paired differences with 10{,}000 bootstrap replicates to obtain a 95\%
confidence interval. We also test the paired null hypothesis with a sign-flip
permutation test that uses 10{,}000 relabellings of the same per-image
differences. We use a permutation test because the per-image counting errors
are small integers with many ties. The tie correction for a signed-rank
statistic is approximate in this setting, whereas the sign-flip null is exact
under the paired-exchangeability hypothesis. Detection recall is a ratio of
pooled counts, not a per-image mean. Its interval therefore uses the same
bootstrap applied to the per-image matched and ground-truth counts. When we test
a family of levels together, we adjust the $p$-values across that family with
the Holm correction.

We average each per-image series over seeds before the paired test. The
resulting interval therefore refers to a recipe and not to one checkpoint.

We report both estimates. The bootstrap interval quantifies uncertainty on this
validation set. The seed range shows how much of a difference remains after
retraining.
